\documentclass[11pt]{article}

\usepackage[letterpaper, top=1in, bottom=1in, left=1.25in, right=1.25in]{geometry}
\usepackage{lmodern}
\usepackage[T1]{fontenc}
\usepackage{amsmath}
\usepackage{amssymb}
\usepackage{graphicx}
\graphicspath{{../outputs/}}
\usepackage{xcolor}
\usepackage{booktabs}
\usepackage{hyperref}
\usepackage{natbib}


% -------------------------------------------------------
% TITLE
% -------------------------------------------------------
\title{\textbf{Temporal Superposition in GPT-2: \\
How Historical Time is Encoded and \\
Distributed Across Layers}}

\author{
  Sam Brown \quad Peyton Calvert \\[0.5em]
  \textit{Your Institution} \\
  \texttt{\{sbrown, pcalvert\}@yourinstitution.edu}
}

\date{}

% -------------------------------------------------------
\begin{document}
% -------------------------------------------------------

\maketitle

% -------------------------------------------------------
\begin{abstract}
% -------------------------------------------------------
% 150-250 words.
% Cover: (1) what you did, (2) how you did it, 
% (3) what you found, (4) why it matters.
% Write this last once the paper is finished.

We investigate how GPT-2-small encodes historical time in its residual 
stream using sparse autoencoders (SAEs). TODO: fill in findings summary.

\end{abstract}

% -------------------------------------------------------
\section{Introduction}
% -------------------------------------------------------
% ~1 page.
% - Motivate the question: why does it matter how LLMs encode time?
% - Brief summary of what you did and what you found
% - Roadmap of the paper ("The rest of this paper is structured as follows...")

Large language models have demonstrated their ability to "know" historical facts, but we have limited understanding of how that
knowledge is stored internally. Understanding the geometric structure of factual knowledge in LLMs is a foundational
question for mechanistic interpretability. 

% -------------------------------------------------------
\section{Related Work}
% -------------------------------------------------------
% ~1 page.
% Cover these papers and how your work relates to each:
% - Gurnee & Tegmark 2023 (linear temporal axis)
% - Engels et al. 2024 (circular features, days/months)
% - Elhage et al. 2022 (superposition hypothesis)
% - Cunningham et al. / Bloom SAEs (SAE methodology)
% - Lei & Cooper 2025 (layerwise knowledge geometry)
% - Calendar feature geometry post (Leask et al.)

TODO

% -------------------------------------------------------
\section{Background}
% -------------------------------------------------------
% ~0.5 page.
% Brief technical background a reader needs before the methods:
% - What is the residual stream in a transformer?
% - What are sparse autoencoders and what do they decompose?
% - What is superposition?
% Keep this concise — assume ML-literate reader.

\subsection{Residual Stream}
Transformer models process tokens by passing representations through a sequence of attention and MLP layers. 
These layers write incremental updates into a shared vector called the residual stream. Formally if $\mathbf{x}^{(\ell)} \in \mathbb{R}^d$
denotes the residual stream at layer $\ell$, each layer computes $\mathbf{x}^{(\ell+1)} = \mathbf{x}^{(\ell)} + f^{(\ell)}$, where 
$f^{(\ell)}$ is the combined attention and MLP transformation. This additive structure means the residual stream accumulates information 
across depth, and that respresentations at any layer can be read as a type of superposition of all contributions made so far.
In this work we probe the residual stream of GPT-2 Small \cite{radford2019}, a 117 million parameter transformer with 12 layers, 
12 attention heads, and a hidden dimension of d = 768. 
\subsection{Superposition}
A naive expectation is that each dimension of the residual stream encodes one interpretable feature. In practice, the number of concepts
a model must represent exceeds its hidden dimension $d$ ($\mathbb{R}^{768}$ in this paper). As a result, models exploit the fact that high 
dimensional spaces can accommodate many nearly-orthogonal directions. This phenomenon is called superposition and it allows a model to 
store a large amount of features at the same time. Features in superposition are distributed across many neurons. 

\subsection{Sparse Autoencoders}
Sparse autoencoders (SAEs) offer a way to recover the underlying precise features from an ambiguous residual stream \cite{cunningham2024}.
An SAE consists of an encoder that projects the d-dimensional activation into a much large latent space $\mathbb{R}^{d_{\text{SAE}}}$
with $d_{\text{SAE}} \gg d$, followed by a ReLU that enforces sparsity and a linear decoder that reconstructs the original activation. 
Training minimises reconstruction error subject to an $L_1$ penalty on the latent activations, encouraging each input to be explained by 
only a small number of active features at a time. The resulting latent dimensions have been shown to correspond to interpretable, monosemantic concepts
\cite{cunningham2024, bloom2024}. In this work we use the SAE of \citet{bloom2024} trained on GPT-2 Small's residual stream. 
% -------------------------------------------------------
\section{Method}
% -------------------------------------------------------
% ~1.5 pages.
% Describe your pipeline clearly enough that someone could reproduce it.

\subsection{Model and SAE}
% GPT-2-small, 12 layers, 768-dim residual stream.
% Joseph Bloom SAEs (gpt2-small-res-jb), 49152 features.
% Hook point: blocks.{layer}.hook_resid_pre

TODO

\subsection{Prompt Design}
% Describe clean_year_items — why no BC/AD tokens.
% Explain the contamination discovery and why you switched.
% Briefly describe the days/months prompts for phase 5.

TODO

\subsection{Temporal Axis Measurement}
% PCA on SAE reconstruction.
% PC1-year Pearson correlation |r| as the metric.
% Why PC1 captures temporal ordering.

TODO

\subsection{Superposition Depth Scan}
% The ablation loop — zeroing one SAE feature at a time.
% Max single-feature drop as the concentration metric.
% Why declining max drop with flat |r| indicates superposition.

TODO

\subsection{Circular Encoding Measurement}
% For days of week and months of year.
% Circle fitting on PC2/PC3 (fit_circle_algebraic).
% Angular correlation as the metric instead of PC1-year r.
% Why linear correlation fails for cyclic concepts.

TODO

% -------------------------------------------------------
\section{Experiments and Results}
% -------------------------------------------------------
% This is the core of the paper. ~3-4 pages including figures.

\subsection{BC/AD Token Contamination}
% Your most methodologically important finding.
% Show phase3 vs phase4 comparison.
% Explain the British Columbia detector (Neuronpedia screenshot).
% Original prompts: |r|~0.95 at layer 0. Clean prompts: |r|~0.40.

TODO

\begin{figure}[h]
  \centering
  \includegraphics[width=0.9\linewidth]{phase4_clean_superposition_depth.png}
  \caption{
    Left: concentration of temporal signal (max single-feature SAE 
    ablation drop) across all 12 layers for clean prompts. 
    Right: baseline temporal axis strength (|PC1-year r|) across layers.
    The temporal axis builds progressively from near zero at layer 0 
    to $|r| \approx 0.93$ by layer 4, then remains stable.
  }
  \label{fig:clean_depth}
\end{figure}

\subsection{Temporal Axis Construction}
% Clean prompts show |r| building from 0.40 to 0.93 across layers 0-4.
% This was hidden in the contaminated dataset.
% Contrast with original prompts where |r| is flat from layer 0.

TODO

\begin{figure}[h]
  \centering
  \includegraphics[width=0.9\linewidth]{phase1_year_linearity.png}
  \caption{
    PC1 vs year value for all historical prompts at layer 4, 
    showing strong linear decodability ($r = -0.920$).
  }
  \label{fig:year_linearity}
\end{figure}

\subsection{Superposition Transition}
% Peak concentration at layer 3 (max drop = 0.285, feature 19325).
% Monotonic decline to near zero by layer 9.
% Baseline |r| stays flat throughout.
% Neuronpedia: feature 19325 is a 19th century year detector.
% Interpretation: construction then progressive superposition.

TODO

\begin{figure}[h]
  \centering
  \includegraphics[width=0.9\linewidth]{phase3_superposition_depth.png}
  \caption{
    Superposition depth scan on original dataset across all 12 layers.
    Left: max single-feature drop. Right: baseline $|r|$.
    Note early layers dominated by BC token detector (feature 22992).
  }
  \label{fig:orig_depth}
\end{figure}

\subsection{Load-Bearing Feature Identification}
% Feature 19325 at layer 3 — Neuronpedia lookup.
% Fires on two-digit suffixes of 1800s year expressions.
% Auto-label: "years in the 1800s" / "historical years".
% Contrast with feature 22992 at layer 1 — British Columbia.
% This is your mechanistic finding.

TODO

\subsection{Circular Concepts: Days of Week and Months of Year}
% Phase 5 results.
% Different pattern from linear time — oscillating concentration.
% Days of week collapse at layers 10-11.
% Comparison plot of all four concepts.

TODO

\begin{figure}[h]
  \centering
  \includegraphics[width=0.9\linewidth]{phase5_circular_superposition_depth.png}
  \caption{
    Superposition depth scan for days of the week and months of the 
    year using angular correlation as the signal metric. Unlike linear 
    temporal encoding, circular concepts show oscillating concentration 
    with no monotonic decline.
  }
  \label{fig:circular_depth}
\end{figure}

\subsection{Days of the Month: Novel Geometric Analysis}
% Phase 6 results.
% Do days of the month show circular geometry?
% Compare angular_r against days of week and months of year.
% This is your novel geometric finding.

TODO

\begin{figure}[h]
  \centering
  \includegraphics[width=0.9\linewidth]{phase6_cyclic_concepts_comparison.png}
  \caption{
    PC2 vs PC3 geometric structure for three cyclic temporal concepts 
    in GPT-2-small. Left: days of the week (layer 7). 
    Middle: months of the year (layer 7). 
    Right: days of the month (best layer). 
    Points colored by position in cycle using HSV colormap.
    Dashed circle shows algebraic least-squares fit.
  }
  \label{fig:cyclic_comparison}
\end{figure}

% -------------------------------------------------------
\section{Discussion}
% -------------------------------------------------------
% ~1 page.
% Interpret your findings. What do they mean?

\subsection{Linear vs Circular Temporal Encoding}
% Why does linear time show monotonic superposition but 
% circular time shows oscillation?
% Hypothesis: circular encoding requires coordinated multi-layer 
% activity, linear encoding can be handed off between layers.

TODO

\subsection{The BC/AD Contamination Finding}
% Methodological contribution.
% Surface tokens can dominate early-layer representations.
% Implication for probing studies generally.
% The shortcut hypothesis — does surface token presence 
% delay genuine semantic encoding?

TODO

\subsection{Limitations}
% Small prompt set (63 prompts for temporal).
% Single model (GPT-2-small).
% SAE reconstruction error (~7% of variance).
% Clustering algorithm optimized for temporal, less suited 
% for geographic/numeric concepts.
% Only one load-bearing feature identified via Neuronpedia.

TODO

% -------------------------------------------------------
\section{Conclusion}
% -------------------------------------------------------
% ~0.5 page.
% Summarize your two main findings:
% 1. Temporal axis lifecycle: construction (layers 0-4), 
%    peak concentration (layer 3), superposition (layers 4-9)
% 2. Linear vs circular encoding behave fundamentally differently
% State what future work should do next.

TODO

% -------------------------------------------------------
\section*{Acknowledgements and Contributions}
% -------------------------------------------------------

\textbf{Sam Brown:} TODO (e.g. data pipeline, phase 1-3 analysis, writing)

\textbf{Peyton Calvert:} TODO (e.g. geometric analysis, phase 5-6, writing)

% -------------------------------------------------------
\section*{Questions for the Instructor}
% -------------------------------------------------------
% Required section per the assignment.
% List any conceptual or methodological questions you have.

\begin{enumerate}
  \item TODO
  \item TODO
\end{enumerate}

% -------------------------------------------------------
% BIBLIOGRAPHY
% -------------------------------------------------------
\bibliographystyle{abbrvnat}
\bibliography{references}

\appendix

% -------------------------------------------------------
\section{Additional Figures}
% -------------------------------------------------------

\begin{figure}[h]
  \centering
  \includegraphics[width=0.9\linewidth]{phase1_all_pca.png}
  \caption{
    PCA of raw residual stream activations at layer 4 for all 
    historical prompts. BC years (dark purple) and modern years 
    (red/orange) separate clearly along PC1 and PC2.
  }
  \label{fig:all_pca}
\end{figure}

\begin{figure}[h]
  \centering
  \includegraphics[width=0.9\linewidth]{phase2_geometric_layer4.png}
  \caption{
    Phase 2 geometric analysis at layer 4. Left to right: raw PCA, 
    PC2 vs PC3 with fitted circle (RMSE=3.235), detrended PCA, 
    angular position vs year ($r=-0.625$). The high circle RMSE 
    indicates the helical/circular hypothesis is not supported 
    for historical time.
  }
  \label{fig:geometric}
\end{figure}

\begin{figure}[h]
  \centering
  \includegraphics[width=0.9\linewidth]{phase1_ablation_histogram.png}
  \caption{
    Distribution of single-feature ablation drops across 60 cluster 
    features at layer 4. One feature causes a drop of $\sim$0.25, 
    while the remaining 59 cause drops near zero.
  }
  \label{fig:ablation_hist}
\end{figure}

\end{document}